% 张三慧V2第23章章末习题23.1--23.19候选独立解析。
% 教材扫描PDF287--289，辅导S PDF321--325；配套公式的疑点单列保留。
\section*{第23章\quad 光的衍射：独立解析}
\paragraph{习题 23.1}
单缝第一零点满足 \(a\sin\theta_1=\lambda\)。
傍轴下焦平面中央明纹\emph{全宽}
\[
\boxed{w=2f\tan\theta_1\simeq\frac{2f\lambda}{a}
=\frac{2(0.5)(546.1\times10^{-9})}{10^{-4}}
=5.461\,\mathrm{mm}}.
\]
\paragraph{习题 23.2}
由第一极小 \(a\sin5^\circ=\lambda\) 得
\[
\boxed{a=\frac{632.8\,\mathrm{nm}}{\sin5^\circ}
\simeq7.26\,\mathrm{\mu m}}.
\]
\paragraph{习题 23.3}
教材近似用第 \(k\) 个次级明纹
\(a\sin\theta_k\simeq(k+\frac12)\lambda\)。
同一方向上
\(\frac72\lambda_x=\frac52(600\,\mathrm{nm})\)，故
\[
\boxed{\lambda_x\simeq\frac57(600)\,\mathrm{nm}
\simeq429\,\mathrm{nm}}.
\]
不应把两个级次颠倒得到840nm。严格的单缝
次级极大须解 \(\tan\beta=\beta\)，会使近似值略变。
\paragraph{习题 23.4}
港口是单一宽缝 \(a=50\,\mathrm m\)；
中央波束在两侧第一极小之间，满足
\[
\boxed{\Delta\theta=2\arcsin\frac{\lambda}{a}
=2\arcsin(20/50)\simeq47.16^\circ}.
\]
不能套圆孔瑞利判据的额外系数1.22。
\paragraph{习题 23.5}
瞳孔圆孔的第一暗环半角
\(\theta_1\simeq1.22\lambda/D\)，故
\[
\boxed{\Delta\theta\simeq2.44\lambda/D
=1.917\times10^{-4}\,\mathrm{rad}}.
\]
视网膜的艾里斑直径
\[
\boxed{d_{\rm img}\simeq l\Delta\theta
=23\,\mathrm{mm}\times1.917\times10^{-4}
\simeq4.41\,\mathrm{\mu m}}.
\]
星像覆盖感光细胞数的面密度估计为
\[
\boxed{N\simeq\frac{\pi d_{\rm img}^2}{4}
(1.5\times10^5\,\mathrm{mm^{-2}})
\simeq2.29},
\]
即2--3个细胞量级。注意不能将半径误当直径：
覆盖面积会错四倍。
\paragraph{习题 23.6}
卫星微波圆孔直径 \(D_1=1.5\times10^3\,\mathrm m\)：
\[
\boxed{\delta=2\arcsin(1.22\lambda/D_1)
\simeq1.63\times10^{-4}\,\mathrm{rad}}.
\]
按配套S PDF322的近轴几何，
考虑光束在起点就具有口径 \(D_1\)，
距离 \(H=3.5\times10^7\,\mathrm m\) 的地面覆盖直径
\[
\boxed{D_{\rm ground}\simeq D_1+2H\tan(\delta/2)
\simeq7.19\,\mathrm{km}}.
\]
若只算远场\emph{发散增加量}，所得约5.69km，
与包含发射口径的7.19km是不同定义；
地球曲率和大气未计。
\paragraph{习题 23.7}
眼睛刚分辨时
\(d_{\rm lights}/L=\theta_R=1.22\lambda/D\)，所以
\[
\boxed{L_{\max}\simeq\frac{d_{\rm lights}D}{1.22\lambda}
=\frac{1.20(.005)}{1.22(550\times10^{-9})}
\simeq8.94\,\mathrm{km}}.
\]
真实眩光、大气湍流可能显著降低此距离。
\paragraph{习题 23.8}
要求识别的角距为
\[
\boxed{\theta\simeq\frac{s}{H}
=\frac{.05}{160000}
=3.125\times10^{-7}\,\mathrm{rad}},
\]
由瑞利判据
\[
\boxed{D_{\min}=\frac{1.22\lambda}{\theta}
=\frac{1.22(500\times10^{-9})}{3.125\times10^{-7}}
\simeq1.95\,\mathrm m}.
\]
这仅是理想镜头的衍射下限，非相机最终设计规格。
\paragraph{习题 23.9}
FAST有效口径是 \(D=300\,\mathrm m\)
（不能用整个500m球面口径），故
\[
\boxed{\theta_R=1.22\frac{.20}{300}
=8.133\times10^{-4}\,\mathrm{rad}
\simeq2.80'}.
\]
\paragraph{习题 23.10}
双望远镜构成的甚长基线 \(B=10^7\,\mathrm m\)，
按合成孔径的相位分辨标度
\[
\boxed{\theta\sim\lambda/B}.
\]
题面只说厘米波段，\emph{无唯一波长}。
若按配套书示例取 \(\lambda=5\,\mathrm{cm}\)，
\[
\boxed{\theta\sim5.0\times10^{-9}\,\mathrm{rad}
\simeq1.03\times10^{-3}\,\text{角秒}}.
\]
不同厘米波长需按比例换算；尚需考虑基线投影、
时钟相位和实际信噪比。
\paragraph{习题 23.11}
双星角距 \(\theta=14''=14\pi/(180\times3600)\,\mathrm{rad}\)，
\[
\boxed{D_{\min}=\frac{1.22\lambda}{\theta}
=\frac{1.22(550\times10^{-9})}{14\pi/(180\times3600)}
\simeq9.89\,\mathrm{mm}\approx1.0\,\mathrm{cm}}.
\]
\paragraph{习题 23.12}
双缝干涉明纹间距
\[
\boxed{\Delta x=f\lambda/d
=0.5(480\times10^{-9})/(10^{-4})=2.4\,\mathrm{mm}}.
\]
单缝包络的中央明纹全宽
\[
\boxed{W=2f\lambda/a
=2(.5)(480\times10^{-9})/(.02\times10^{-3})
=24\,\mathrm{mm}}.
\]
干涉级次 \(m\) 落在中央包络内必须
\(|m|<d/a=5\)。
第 \(m=\pm5\) 级正位于单缝暗纹而缺级，
允许 \(m=-4,-3,\dots,4\)，总数
\(\boxed{9\text{个干涉主极大}}\)。
\paragraph{习题 23.13}
栅距 \(d=W/N=.020/6000
=3.333\times10^{-6}\,\mathrm m\)。
主极大满足 \(d\sin\theta_m=m\lambda\)，所以
\[
\boxed{\sin\theta_m=0.17679m,\quad m=0,\pm1,\ldots,\pm5}.
\]
角位置约为
\[
\boxed{\theta=0^\circ,\ \pm10.18^\circ,\ \pm20.71^\circ,\
\pm32.04^\circ,\ \pm44.99^\circ,\ \pm62.13^\circ}.
\]
第6级所需 \(|\sin\theta|>1\)，不存在。
\paragraph{习题 23.14}
光栅常量 \(d=0.01/6000=1.667\,\mathrm{\mu m}\)；
\[
\boxed{\lambda=d\sin20^\circ\simeq570\,\mathrm{nm}},
\qquad
\boxed{\theta_2=\arcsin(2\sin20^\circ)
\simeq43.2^\circ}.
\]
另一侧的第2级同样位于负角度。
\paragraph{习题 23.15}
原图23.23的 \(H_\alpha\) 谱线第2级
\(k=2,\theta\simeq41^\circ,\lambda=656.3\,\mathrm{nm}\)，
故
\[
\boxed{d\simeq\frac{2\lambda}{\sin41^\circ}
\simeq2.00\,\mathrm{\mu m}}.
\]
图上第3级缺级，即第3级光栅主极大与单缝
第一极小重合：\(3\lambda/d=\lambda/a\)，
\[
\boxed{a\simeq d/3\simeq0.667\,\mathrm{\mu m}}.
\]
基于图示约41°读数，因此是估计值而非实测真值。
\paragraph{习题 23.16}
分辨本领 \(R=kN=\lambda/\Delta\lambda\)。
第1级 \(k=1\)，需整数
\[
N\ge656.3/0.18=3646.111\ldots
\quad\Longrightarrow\quad\boxed{N_{\min}=3647}.
\]
配套S PDF324写3646，按严格最小整数
应\emph{向上取整}而非截断；
报告数量级可作约 \(3.65\times10^3\) 条。
\paragraph{习题 23.17}
波长 \(\lambda=c/f\simeq1.292\,\mathrm m\)；
同相的一维 \(N=192\) 阵列因子
\[
I(\theta)\propto
\left[\frac{\sin(N\beta)}{\sin\beta}\right]^2,\qquad
\beta=\frac{\pi d\sin\theta}{\lambda},\quad d=6\,\mathrm m.
\]
正天顶 \(\theta=0\) 时为\emph{主极大}。
首次零点 \(N\beta=\pi\)：
\[
\boxed{\theta_{\min,1}=\arcsin\frac{\lambda}{Nd}
\simeq1.122\times10^{-3}\,\mathrm{rad}
\simeq3.86'}.
\]
第一\emph{非零级主极大}为 \(d\sin\theta=\lambda\)，
\[
\boxed{\theta_{\rm grating,1}=\arcsin(\lambda/d)\simeq12.5^\circ}.
\]
严格区分：这不是离天顶最近的\emph{局部极大}。
主瓣与下一级主极大之间尚有一系列低强度旁瓣峰；
第一旁瓣峰约 \(1.4303\lambda/(Nd)
\simeq1.60\times10^{-3}\,\mathrm{rad}\)。
原配套S PDF325报12.4°是\emph{第一非零级主极大}。
\paragraph{习题 23.18}
布拉格式 \(2d\sin\varphi=k\lambda\)，且
\[
2d\sin45^\circ=2(.275)\sin45^\circ
\simeq0.3889\,\mathrm{nm}.
\]
区间 \(0.095\le0.3889/k\le0.130\)
仅允许 \(k=3,4\)，因而
\[
\boxed{\lambda_3\simeq0.1296\,\mathrm{nm},\quad
\lambda_4\simeq0.0972\,\mathrm{nm}}.
\]
\paragraph{习题 23.19}
按教材图23.34相邻\emph{面内}原子
对电子的散射波程差
\(\Delta=d\sin\varphi\)。
在 \(\varphi=50^\circ\) 处一级相长，
\[
\boxed{\lambda=d\sin50^\circ
=0.215\sin50^\circ\,\mathrm{nm}
\simeq0.165\,\mathrm{nm}}.
\]
不能未经几何转换直接将晶面间距 \(d\) 套入
厚度方向的布拉格双程公式。
